%! Orthogonal Matrices and Gram--Schmidt — Unit 4, Lesson 4.4 — Warm-Up
\documentclass[10pt]{article}
\usepackage{linalg-article}
\usepackage{linalg-boxes}
\newcommand{\TallMath}[1]{$\displaystyle #1\rule[-1.4em]{0pt}{3.2em}$}
\newcommand{\vv}{\mathbf{v}}
\newcommand{\ww}{\mathbf{w}}
\newcommand{\xx}{\mathbf{x}}
\newcommand{\bb}{\mathbf{b}}
\newcommand{\T}{\mathsf{T}}

\begin{document}

\pageheader{Unit 4, Lesson 4.4}{Warm-Up}
\namedateperiod

\vspace{0.12in}

\begin{notesbox}{Getting Ready: Skills We'll Use Today}
\small These three ideas (Lessons 1.2, 4.0, 4.3) set up today's lesson. Answer quickly.

\vspace{4pt}
\textbf{1. Length and unit vectors.} For $\vv=(3,4)$:
\[
  \|\vv\|=\sqrt{\vv\cdot\vv}=\blank{1.6cm}\,. \qquad
  \text{A unit vector in the same direction is } \frac{\vv}{\|\vv\|}=\blank{2.4cm}\,.
\]

\vspace{4pt}
\textbf{2. The perpendicular test.} Are $\vv=(3,4)$ and $\ww=(4,-3)$ perpendicular? Compute the dot
product and decide.
\[
  \vv\cdot\ww=\blank{2.0cm} \qquad\Longrightarrow\qquad \text{perpendicular?}\ \blank{2.0cm}
\]

\vspace{4pt}
\textbf{3. Last lesson's normal equations.} To fit data we solved $A^{\T}A\hat{\xx}=A^{\T}\bb$. Fill in: if
$A^{\T}A$ turned out to be the identity $I$, then $\hat{\xx}=\blank{2.4cm}$ --- and there is
\blank{2.4cm} system left to solve.
\end{notesbox}

\end{document}
